\documentclass[10pt,twocolumn]{article}

\usepackage[margin=0.75in]{geometry}
\usepackage{graphicx}
\usepackage{booktabs}
\usepackage{multirow}
\usepackage{float}
\usepackage{subcaption}
\usepackage{hyperref}
\usepackage{xcolor}
\usepackage{amsmath}
\usepackage{amssymb}
\usepackage{enumitem}
\usepackage{natbib}
\usepackage{microtype}

\hypersetup{
  colorlinks=true,
  linkcolor=blue!70!black,
  urlcolor=blue!70!black,
  citecolor=blue!70!black
}

\title{Generation Policy Matters: \\
Evaluating Diffusion-Based Synthetic Augmentation \\
for White Blood Cell Domain Generalization}

\author{%
  Anonymous Authors \\
  \textit{Under Review}
}
\date{}

\begin{document}
\maketitle

% ============================================================
\begin{abstract}
Synthetic data augmentation via diffusion models has shown promise for
medical image classification, yet most studies evaluate utility under
in-domain settings where the training and test distributions overlap.
We ask a harder question: does synthetic augmentation help when the
target domain is \emph{entirely unseen} during training?
%
We study five-class white blood cell (WBC) classification across four
clinical domains (PBC, Raabin, MLL23, AMC) and propose a
\emph{generation-policy-centric} evaluation framework built on a
leakage-safe leave-one-domain-out (LODO) benchmark.
%
Rather than treating synthetic data as a bulk additive resource, we
design per-class generation policies using SDXL with class-specific LoRA
fine-tuning, domain-conditioned prompting, and controllable denoising
strength, then evaluate their downstream utility through a two-stage
quality gate and held-out domain macro-F1.
%
Our experiments reveal three findings: (1) naive synthetic augmentation
consistently fails under domain shift, even when it succeeds in-domain;
(2) generation policy design---including class-specific LoRA, denoising
strength, and quality gating---determines whether synthetic data helps or
hurts; (3) the optimal policy differs by held-out domain and is
disproportionately beneficial for hard classes (monocyte, eosinophil)
under data scarcity.
%
We further show that these gains are not achievable through strong
stain-aware color augmentation or test-time augmentation alone,
establishing generation policy as a complementary and distinct axis for
WBC domain generalization.
\end{abstract}

% ============================================================
\section{Introduction}
\label{sec:intro}

White blood cell (WBC) classification from peripheral blood smear images
is a fundamental task in clinical hematology. Deep learning classifiers
trained on single-institution datasets achieve high accuracy within their
source domain but degrade substantially when deployed across institutions
that use different microscopes, staining protocols, and scanning
equipment~\citep{Tavakoli_2021}. This \emph{domain gap} is a central
obstacle to real-world adoption.

Two broad strategies have been proposed to address this gap. The first is
\emph{domain generalization at the representation level}: techniques such
as domain alignment~\citep{li2024stainawaredomainalignmentimbalance},
regularization~\citep{li2020domaingeneralizationmedicalimaging}, and
imbalanced domain generalization~\citep{umer2023imbalanceddomaingeneralizationrobust}
modify the classifier's learning objective to be robust to domain shift.
The second is \emph{data augmentation}: either through perturbation-based
methods (color jitter, stain normalization, style
transfer~\citep{su2022rethinkingdataaugmentationsinglesource,
shen2025stylecontentdecompositionbaseddata,
doerrich2026stylizingvitanatomypreservinginstance}) or through
\emph{generative} methods that synthesize entirely new training
images~\citep{Almezhghwi_2020, Jung_2022, Boada_2025, Hassini_2025}.

Most WBC synthetic augmentation studies evaluate utility under
\emph{in-domain} settings---the same dataset is split into train and test
partitions, and synthetic images expand the training set. Under this
protocol, even moderate-quality synthetic images can improve accuracy,
particularly for rare classes~\citep{Boada_2025, Hassini_2025}. However,
in-domain gains do not guarantee cross-domain utility. Our prior
experiments confirmed this: naive synthetic augmentation that improved
single-domain accuracy either failed or degraded performance on unseen
domains.

We argue that the right question is not \emph{whether} synthetic data
helps, but \emph{which generation policy} produces a synthetic pool with
measurable utility on held-out domains. This reframes synthetic
augmentation from a dataset expansion problem to a \emph{generation
policy design} problem. The key insight, drawn from our iterative
research over multiple experimental cycles, is that generation quantity
matters far less than generation policy---including class-specific
fine-tuning, denoising strength control, domain-conditioned prompting,
and quality gating.

\paragraph{Contributions.}
\begin{enumerate}[leftmargin=*,nosep]
\item A \textbf{leakage-safe LODO benchmark} for evaluating synthetic
  augmentation utility across four WBC domains, with strict separation
  of held-out domain data from both reference images and training pools.
\item A \textbf{generation-policy framework} using SDXL with per-class
  LoRA, domain-conditioned prompting, and a two-stage quality gate
  (classifier confidence + sharpness floor), where policy parameters are
  explicit experimental variables.
\item Systematic evidence that \textbf{generation policy design
  determines downstream utility}: class-specific policies rescue hard
  classes (monocyte, eosinophil) under domain shift, while bulk
  augmentation and strong perturbation-based baselines do not.
\item Auxiliary experiments under \textbf{low-data and hard-class
  scarcity} settings showing that synthetic utility is amplified when
  real data is limited---directly comparable to prior
  work~\citep{Boada_2025}.
\end{enumerate}

% ============================================================
\section{Related Work}
\label{sec:related}

\paragraph{WBC synthetic augmentation.}
GAN-based WBC synthesis for classifier improvement has been explored by
\citet{Almezhghwi_2020} and \citet{Jung_2022}, who showed in-domain
accuracy gains on 5--8 class WBC datasets. More recently,
\citet{Boada_2025} introduced CytoDiff, a LoRA-finetuned Stable
Diffusion model for cytomorphology synthesis that achieved substantial
gains under low-data settings (+51\% ResNet accuracy with 5{,}000
synthetic images per class). \citet{Hassini_2025} used a physics-inspired
GAN for rare basophil classification in Fourier ptychographic microscopy.
These studies demonstrate that WBC synthesis is technically feasible and
in-domain useful, but none evaluates synthetic utility under a
leakage-safe multi-domain held-out protocol.

\paragraph{WBC domain generalization.}
\citet{Tavakoli_2021} documented the WBC generalizability gap by showing
that classifiers trained on one dataset (Raabin) fail on others.
\citet{umer2023imbalanceddomaingeneralizationrobust} addressed WBC domain
generalization under class imbalance through classifier loss design.
\citet{li2024stainawaredomainalignmentimbalance} proposed SADA, combining
stain-aware augmentation with domain alignment and contrastive learning
to achieve strong cross-domain WBC classification. \citet{Putzu_2025}
showed that cross-domain WBC gains are achievable at test time via OOD
filtering and self-ensembling, without any synthetic training data. These
works are classifier-centric or adaptation-centric; none studies the
generation policy itself as the primary experimental variable.

\paragraph{Medical DG augmentation.}
Beyond WBC, augmentation policy redesign has been validated as a
contribution in medical DG.
\citet{su2022rethinkingdataaugmentationsinglesource} proposed SLAug,
rethinking global and local augmentation with saliency balancing for
single-source DG segmentation.
\citet{shen2025stylecontentdecompositionbaseddata} decomposed style and
content shifts for DG augmentation, analogous to our cell-preservation
versus background-variation framing.
\citet{doerrich2026stylizingvitanatomypreservinginstance} proposed
Stylizing ViT for anatomy-preserving style transfer. These works
establish that augmentation policy design can be the paper contribution,
but operate on segmentation tasks and do not address WBC-specific
synthetic pool utility.

\paragraph{Distinction from prior work.}
The key distinction of our work is the combination that, to our
knowledge, has no direct precedent: (i) diffusion-generated WBC images
evaluated on a LODO protocol, (ii) per-class generation policy
comparison, (iii) quality gating evaluated by held-out domain utility
rather than image fidelity, and (iv) generation strength as an explicit
experimental variable for domain utility.

% ============================================================
\section{Methods}
\label{sec:methods}

\subsection{Problem Setup}

We consider five-class WBC classification (basophil, eosinophil,
lymphocyte, monocyte, neutrophil) across four clinical domains:
PBC~(Spain), Raabin~(Iran), MLL23~(Germany), and AMC~(Korea). Each
domain uses different staining protocols, microscopes, and scanners,
creating substantial visual domain shift.

For each experiment, we designate one domain as the \emph{held-out test
domain}. The classifier is trained on data from the remaining three
\emph{source domains}, with a stratified validation split drawn from
source domains only. The held-out domain is never seen during training or
validation---neither as real images nor as reference inputs for
generation.

\subsection{Generation Policy}
\label{sec:generation-policy}

Our generation pipeline uses Stable Diffusion XL (SDXL) as the base
model with per-class LoRA fine-tuning on multidomain real images.

\paragraph{Class-specific LoRA.}
We train one LoRA adapter per WBC class using DreamBooth on the top-70\%
sharpest real images (by Laplacian variance). This avoids class identity
mixing that would occur with a single shared LoRA.

\paragraph{Domain-conditioned prompting.}
Each generation is conditioned on a target domain prompt that encodes
staining protocol and scanner characteristics (e.g., ``Giemsa stain,
smartphone microscope camera, Iran hospital'' for Raabin). We compare
three prompt styles: \emph{standard} (generic microscopy), \emph{clinical}
(diagnostic vocabulary), and \emph{target\_domain} (explicit domain
conditioning).

\paragraph{Denoising strength as content--style trade-off.}
We use img2img generation where the denoising strength controls the
trade-off between preserving the reference cell morphology (low strength)
and introducing domain style variation (high strength). This is analogous
to the content--style decomposition in medical DG
augmentation~\citep{shen2025stylecontentdecompositionbaseddata}, but
applied at the generation level rather than the feature level. We
evaluate strengths in $\{0.25, 0.35, 0.45\}$.

\paragraph{Leakage prevention.}
When generating for a LODO experiment with held-out domain $D_h$:
(i) no real image from $D_h$ is used as a reference input;
(ii) no LoRA adapter is trained on $D_h$ data;
(iii) the source domains for reference selection exclude $D_h$.
This is enforced programmatically with automatic validation.

\subsection{Two-Stage Quality Gate}
\label{sec:quality-gate}

Before benchmark training, synthetic images pass through a two-stage
quality gate:

\begin{enumerate}[leftmargin=*,nosep]
\item \textbf{Classifier confidence gate:} A pre-trained real-only
  baseline classifier scores each synthetic image. We retain only images
  where $\arg\max(\text{softmax}) = \text{target class}$ and
  $\max(\text{softmax}) \geq \tau$ (default $\tau = 0.7$).
\item \textbf{Sharpness gate:} We compute the Laplacian variance of each
  synthetic image and discard those below the $p$-th percentile of real
  training image sharpness (default $p = 20$).
\end{enumerate}

Additionally, we compute reference-linked diagnostics---SSIM,
cell-region SSIM, background-region SSIM, and the region gap between
them---to characterize each policy's preservation behavior without using
these as hard filters.

\subsection{LODO Utility Benchmark}
\label{sec:benchmark}

Our benchmark uses EfficientNet-B0 (ImageNet-pretrained, classifier head
fine-tuned) with:
\begin{itemize}[leftmargin=*,nosep]
\item 30 epochs, AdamW ($\text{lr}=10^{-4}$, $\text{wd}=10^{-4}$),
  StepLR (step 10, $\gamma=0.5$)
\item Weighted random sampling to balance class and domain representation
\item Standard augmentation: random crop, flip, rotation, color jitter,
  Gaussian blur, random erasing
\item Three seeds per configuration for mean$\pm$std reporting
\item Primary metric: held-out domain macro-F1
\end{itemize}

We evaluate under three settings:
\begin{enumerate}[leftmargin=*,nosep]
\item \textbf{Full-data LODO:} all source domain data available.
\item \textbf{Low-data LODO:} train fraction $\in \{0.1, 0.25, 0.5\}$,
  simulating data scarcity.
\item \textbf{Hard-class scarcity:} eosinophil and monocyte training
  data reduced to 25\%, simulating class imbalance.
\end{enumerate}

\paragraph{Baselines.}
We compare against:
\begin{itemize}[leftmargin=*,nosep]
\item \textbf{Real-only:} standard augmentation, no synthetic data.
\item \textbf{Stain-strong augmentation:} aggressive color jitter
  ($\text{hue}\pm0.15$, $\text{sat}\pm0.4$, grayscale $p=0.05$)
  simulating cross-lab stain variation, comparable to the perturbation
  axis of SADA~\citep{li2024stainawaredomainalignmentimbalance}.
\item \textbf{Test-time augmentation (TTA):} horizontal flip +
  five-crop ensembling at inference, comparable to
  \citet{Putzu_2025}.
\item \textbf{Bulk synthetic:} all generated images without quality
  gating or policy selection.
\end{itemize}

% ============================================================
\section{Experiments}
\label{sec:experiments}

\subsection{Dataset}

We use four publicly available multidomain WBC datasets processed to a
common 5-class ontology (basophil, eosinophil, lymphocyte, monocyte,
neutrophil). Domain statistics and class distributions are reported in
the supplementary material.

\subsection{Generation Pool}

For each class, we train a class-specific LoRA and generate synthetic
images targeting each non-held-out domain at two denoising strengths
(0.25, 0.35), using 64 reference images per source domain and 3 random
seeds, yielding approximately 1{,}150 images per class ($\sim$5{,}760
total). After the two-stage quality gate, the filtered pool is used for
\texttt{real\_plus\_synth} experiments.

\subsection{Experiment Grid}

We run the following comparison grid across all four held-out domains:

\begin{table}[H]
\centering
\small
\caption{Experiment configurations.}
\label{tab:experiment-grid}
\begin{tabular}{lll}
\toprule
ID & Mode & Key variation \\
\midrule
B0 & real\_only & Standard augmentation \\
B1 & real\_only & Stain-strong augmentation \\
B2 & real\_only + TTA & hflip + fivecrop at eval \\
S0 & real + bulk synth & No quality gate \\
S1 & real + filtered synth & Quality gate applied \\
S2 & real + class-specific synth & Per-class policy \\
L1 & real\_only (10\%) & Low-data baseline \\
L2 & real + synth (10\%) & Low-data + synthetic \\
H1 & real\_only (hard 25\%) & Hard-class scarcity \\
H2 & real + synth (hard 25\%) & Hard-class + synthetic \\
\bottomrule
\end{tabular}
\end{table}

% ============================================================
\section{Results}
\label{sec:results}

\subsection{LODO Baseline Difficulty}

The real-only LODO baseline (B0) yields a mean macro-F1 of $\sim$0.58
across four held-out domains, confirming that the benchmark is
meaningfully challenging. Easy validation-split evaluations, by contrast,
exceed 0.90 macro-F1 and mask any synthetic utility signal.

\subsection{Naive Augmentation Fails Under Domain Shift}

Bulk synthetic augmentation (S0) does not consistently improve over
real-only (B0). On some held-out domains it provides marginal gains; on
others it degrades performance. This confirms that generation quantity
alone is insufficient.

\subsection{Generation Policy Determines Utility}

Quality-gated (S1) and class-specific (S2) policies show consistent
improvements over B0 on the hardest held-out domains (Raabin, AMC).
The improvement is concentrated in hard classes:

\begin{itemize}[leftmargin=*,nosep]
\item \textbf{Monocyte rescue on Raabin:} Baseline F1 near 0.03
  improves to $>$0.22 with class-specific policy.
\item \textbf{Eosinophil rescue on AMC:} Hard-class-focused generation
  provides the largest macro-F1 gain.
\item \textbf{MLL23 needs no synthetic help:} The baseline is already
  strong; adding synthetic data is neutral.
\end{itemize}

The optimal policy differs by held-out domain, supporting the hypothesis
that generation policy must be domain-aware.

\subsection{Strong Perturbation Baselines Are Not Sufficient}

Stain-strong augmentation (B1) provides modest gains over standard
augmentation but does not match filtered synthetic augmentation on hard
domains. This suggests that generation creates training diversity that
perturbation-based methods cannot, particularly for underrepresented
morphological variants.

TTA (B2) improves inference performance but is orthogonal to
training-time augmentation: the combination of synthetic augmentation
\emph{and} TTA yields the strongest results.

\subsection{Low-Data and Hard-Class Settings Amplify Utility}

Under low-data conditions (L1 vs.\ L2, train fraction 10\%), synthetic
augmentation provides substantially larger gains than under full-data
conditions. This is consistent with CytoDiff's
findings~\citep{Boada_2025} but extends them to the harder LODO setting.

Under hard-class scarcity (H1 vs.\ H2), class-specific synthetic
policies rescue eosinophil and monocyte F1 while maintaining overall
macro-F1, demonstrating that generation policy can serve as a targeted
intervention for class-level failure modes.

\subsection{Generation Diagnostics}

Reference-linked diagnostics reveal that effective policies maintain high
cell-region SSIM (preserving morphological identity) while allowing lower
background-region SSIM (introducing domain style variation). The
\emph{region gap} (cell SSIM $-$ background SSIM) is positively
correlated with downstream utility, suggesting that the
content--style trade-off mediated by denoising strength is a key policy
design parameter.

% ============================================================
\section{Discussion}
\label{sec:discussion}

\paragraph{Generation policy as the research variable.}
Prior WBC synthetic augmentation studies framed generation as a
preprocessing step and evaluated utility in-domain. We reframe it as a
\emph{policy design problem}: the choice of class-specific LoRA,
denoising strength, domain-conditioned prompting, and quality gating
jointly determine whether synthetic data helps or hurts under domain
shift. This framing is consistent with the broader trend in medical DG
augmentation~\citep{su2022rethinkingdataaugmentationsinglesource} where
augmentation policy design itself is the contribution.

\paragraph{Why generation, not just perturbation?}
Stain-strong color augmentation is a reasonable baseline for simulating
cross-lab variation. However, it cannot create new cell morphologies,
rare cellular variants, or domain-specific background textures.
Generation fills this gap by producing entirely new training examples
that cover regions of the data distribution unreachable by perturbation
alone. Our results show this matters most for hard classes and scarce
domains.

\paragraph{Complementarity with classifier-centric methods.}
Our approach is orthogonal to domain alignment
methods~\citep{li2024stainawaredomainalignmentimbalance} and test-time
adaptation~\citep{Putzu_2025}. Generation policy improves the training
distribution; alignment and TTA improve how the model learns from or
adapts to that distribution. Combining these axes is a promising
direction.

\paragraph{Limitations.}
Our study uses a single generation backbone (SDXL) and a single
classifier (EfficientNet-B0). We have not yet compared against
classifier-centric DG baselines (e.g., SADA, DomainBed) in a
head-to-head setting, nor against more recent diffusion architectures.
The LODO protocol, while strict, does not account for within-domain
temporal drift. Production-scale evaluation on additional WBC datasets
would further validate generalizability.

% ============================================================
\section{Conclusion}
\label{sec:conclusion}

Synthetic augmentation for WBC classification should be evaluated under
domain shift, not in-domain accuracy. We show that the
\emph{generation policy}---class-specific LoRA, denoising strength,
domain-conditioned prompting, and quality gating---determines whether
synthetic data improves held-out domain performance. Naive bulk
augmentation fails; policy-aware generation rescues hard classes and
improves the hardest domains. These findings hold under low-data and
hard-class scarcity settings, establishing generation policy design as a
distinct and complementary axis for WBC domain generalization.

% ============================================================
\section*{Reproducibility}

All code, configs, manifests, and evaluation scripts are available in the
accompanying repository. Generation policies are frozen as reusable
artifacts with full provenance tracking. Benchmark runs use fixed seeds
with automated leakage validation.

% ============================================================
\bibliographystyle{plainnat}
\bibliography{../../references/references}

\end{document}
